\documentclass[conference]{IEEEtran}

\usepackage[utf8]{inputenc}
\usepackage[T1]{fontenc}
\usepackage[spanish,es-nodecimaldot]{babel}
\usepackage{amsmath}
\usepackage{booktabs}
\usepackage{graphicx}
\usepackage{hyperref}
\usepackage{listings}
\usepackage{xcolor}

\hypersetup{
    colorlinks=true,
    linkcolor=blue,
    urlcolor=blue
}

\definecolor{codegray}{RGB}{245,245,245}
\lstset{
    basicstyle=\ttfamily\scriptsize,
    backgroundcolor=\color{codegray},
    frame=single,
    breaklines=true,
    showstringspaces=false
}

\title{Forecasting de Temperatura con RNN y LSTM\\
       usando MLflow, DVC y Docker}

\author{
\IEEEauthorblockN{Jose Miguel Fonseca Cordoba}
\IEEEauthorblockA{Curso de Deep Learning\\
2026}
}

\begin{document}
\maketitle

\begin{abstract}
Se desarrolló un sistema reproducible de forecasting de temperatura usando
el conjunto de datos Jena Climate. Se implementaron y compararon una red
neuronal recurrente (RNN) y una red Long Short-Term Memory (LSTM). El
preprocesamiento utiliza normalización calculada únicamente sobre el conjunto
de entrenamiento y ventanas de 24 observaciones horarias para predecir la
temperatura siguiente. Los experimentos se registraron con MLflow, el
entrenamiento se ejecutó dentro de Docker con GPU y el mejor modelo se
desplegó mediante FastAPI en un contenedor separado. La configuración
experimental mejorada obtuvo un RMSE de $0.7253\,^\circ$C con LSTM, frente a
$0.7330\,^\circ$C$ con RNN. El pipeline también incorpora DVC local,
pre-commit, GitHub Actions y pruebas de inferencia mediante curl.
\end{abstract}

\begin{IEEEkeywords}
series temporales, forecasting, RNN, LSTM, MLflow, DVC, Docker.
\end{IEEEkeywords}

\section{Introducción}

El objetivo del proyecto es construir un pipeline completo de aprendizaje
profundo para pronosticar una serie temporal y desplegar el modelo ganador
como un servicio de inferencia. El sistema debía incluir preprocesamiento,
comparación entre RNN y LSTM, seguimiento experimental con MLflow,
entrenamiento en contenedores, versionado de datos y un servidor de
inferencia sin utilizar Triton.

La solución se implementó como una nueva fase independiente dentro del
repositorio existente, sin modificar las fases anteriores de clasificación
con VGG11 ni detección de objetos.

\section{Datos y preprocesamiento}

Se utilizó Jena Climate, un registro multivariable de observaciones
meteorológicas muestreadas aproximadamente cada diez minutos. Para este
proyecto se seleccionó la variable de temperatura y se tomó una observación
cada seis registros, obteniendo una frecuencia aproximada de una hora.

Los datos se dividieron temporalmente en 70\% para entrenamiento, 15\% para
validación y 15\% para prueba. La media y la desviación estándar se
calcularon exclusivamente con el conjunto de entrenamiento:

\[
x_{\mathrm{norm}} = \frac{x-\mu_{\mathrm{train}}}
                         {\sigma_{\mathrm{train}}}
\]

Cada muestra contiene 24 valores consecutivos y el objetivo es el siguiente
valor:

\[
[x_t,\ldots,x_{t+23}] \longrightarrow x_{t+24}
\]

El preprocesamiento evita fuga de información y se implementa en
\texttt{forecasting/data/series.py}.

\section{Modelos}

Se entrenaron dos modelos con la misma tarea y los mismos datos.

\subsection{RNN}

La RNN utiliza una capa recurrente con una representación oculta de 64
unidades en la configuración experimental. La salida correspondiente al
último instante temporal se pasa a una capa lineal que produce una
predicción escalar.

\subsection{LSTM}

La LSTM utiliza la misma dimensión oculta y una capa lineal de salida. Sus
estados de memoria y compuertas permiten conservar dependencias temporales
de mayor duración, lo cual resulta útil en series con patrones diarios y
estacionales.

\section{Entrenamiento y seguimiento}

El entrenamiento se ejecuta dentro de Docker utilizando una imagen de
PyTorch con CUDA. La configuración mejorada utiliza:

\begin{itemize}
    \item 40 épocas como máximo.
    \item \texttt{hidden\_size=64}.
    \item Tasa de aprendizaje $5\times10^{-4}$.
    \item Batch de 64 muestras.
    \item Early stopping con paciencia de 5 épocas.
    \item Semilla aleatoria fija.
\end{itemize}

MLflow registra los parámetros, las pérdidas por época, la mejor pérdida de
validación y los modelos serializados. El experimento mejorado se identificó
como:

\begin{lstlisting}
jena-temperature-forecasting-experiments
\end{lstlisting}

El entrenamiento se puede reproducir mediante:

\begin{lstlisting}[language=bash]
dvc repro
\end{lstlisting}

DVC mantiene la definición del pipeline en \texttt{dvc.yaml} y el estado
reproducible en \texttt{dvc.lock}. El almacenamiento DVC utilizado en esta
versión es local.

\section{Resultados}

La configuración inicial, con 10 épocas y 32 unidades ocultas, obtuvo:

\begin{table}[h]
\centering
\caption{Resultados de la configuración baseline.}
\begin{tabular}{lcc}
\toprule
Modelo & MSE normalizado & RMSE ($^\circ$C) \\
\midrule
LSTM & 0.007481 & 0.7486 \\
RNN  & 0.007649 & 0.7570 \\
\bottomrule
\end{tabular}
\end{table}

La configuración experimental mejorada obtuvo:

\begin{table}[h]
\centering
\caption{Resultados de la configuración mejorada.}
\begin{tabular}{lcc}
\toprule
Modelo & MSE normalizado & RMSE ($^\circ$C) \\
\midrule
LSTM & 0.007023 & \textbf{0.7253} \\
RNN  & 0.007172 & 0.7330 \\
\bottomrule
\end{tabular}
\end{table}

La LSTM mejorada redujo el RMSE respecto a la LSTM baseline:

\[
0.7486 - 0.7253 = 0.0233\,^\circ\mathrm{C}
\]

Por lo tanto, la LSTM fue seleccionada como modelo final.

\section{Servidor de inferencia}

El modelo final se copia en:

\begin{lstlisting}
forecasting/artifacts/best_model.pt
\end{lstlisting}

El servidor FastAPI se ejecuta en Docker y expone el puerto 5002. El
endpoint de salud confirmó:

\begin{lstlisting}
{
  "status": "ok",
  "model": "lstm",
  "device": "cpu",
  "sequence_length": 24
}
\end{lstlisting}

El endpoint \texttt{POST /predict} recibe exactamente 24 temperaturas y
devuelve la predicción desnormalizada en grados Celsius. El entrenamiento
utiliza GPU, mientras que el servicio de inferencia final utiliza CPU.

\section{Pruebas funcionales}

Se preparó un script con tres pruebas reproducibles:

\begin{enumerate}
    \item Verificación de \texttt{GET /health}.
    \item Predicción válida con 24 valores.
    \item Rechazo de una secuencia inválida con HTTP 422.
\end{enumerate}

El script se encuentra en:

\begin{lstlisting}
forecasting/tests/curl_tests.sh
\end{lstlisting}

La ejecución produjo:

\begin{lstlisting}
All inference curl tests passed.
\end{lstlisting}

Además, las pruebas unitarias finalizaron con:

\begin{lstlisting}
2 passed
\end{lstlisting}

\section{Reproducibilidad y calidad}

El proyecto utiliza Git para el código, DVC para el pipeline y artefactos
grandes, MLflow para el seguimiento de experimentos, Docker para el
entrenamiento e inferencia, y pre-commit para validaciones locales.

GitHub Actions ejecuta automáticamente:

\begin{itemize}
    \item Pruebas unitarias.
    \item Compilación de los módulos Python.
    \item Hooks de pre-commit.
\end{itemize}

El workflow no entrena modelos en cada commit porque el entrenamiento con GPU
no es apropiado para una validación rápida de código.

\section{Limitaciones}

El modelo utiliza únicamente temperatura como variable de entrada y predice
un solo paso temporal. El servidor no realiza todavía forecasting
multihorizonte. Asimismo, el almacenamiento DVC es local, por lo que otra
computadora necesitaría una copia del dataset y de los checkpoints o un
remoto DVC configurado.

\section{Conclusiones}

Se implementó un pipeline completo de forecasting con RNN y LSTM. La
configuración experimental con mayor capacidad y early stopping mejoró el
RMSE de la LSTM de 0.7486 a 0.7253 grados Celsius. El modelo final se
registró con MLflow, se seleccionó mediante el pipeline DVC y se desplegó
correctamente en un servidor FastAPI dentro de Docker.

Las tres pruebas de inferencia y las pruebas unitarias confirmaron el
funcionamiento del servicio. El resultado cumple los requisitos de
preprocesamiento, comparación de modelos, seguimiento experimental,
entrenamiento en contenedor, versionado y despliegue.

\begin{thebibliography}{00}

\bibitem{jena}
M. L. B. et al., ``The Jena Climate dataset,'' Max Planck Institute for
Biogeochemistry, 2009--2016.

\bibitem{mlflow}
MLflow, ``An open source platform for the machine learning lifecycle,''
\url{https://mlflow.org/}.

\bibitem{dvc}
DVC, ``Data Version Control,'' \url{https://dvc.org/}.

\bibitem{lstm}
S. Hochreiter and J. Schmidhuber, ``Long short-term memory,''
\textit{Neural Computation}, vol. 9, no. 8, pp. 1735--1780, 1997.

\end{thebibliography}

\end{document}
